\documentclass[12pt,a4paper]{article}
\usepackage{fontspec}
\IfFontExistsTF{Arial}{\setsansfont{Arial}}{}
\IfFontExistsTF{Consolas}{\setmonofont{Consolas}[Scale=0.9]}{}
\IfFontExistsTF{Times New Roman}{\setmainfont{Times New Roman}}{}
\usepackage[brazil]{babel}
\usepackage[margin=2.6cm]{geometry}
\usepackage{booktabs,tabularx}
\usepackage[hidelinks]{hyperref}
\hypersetup{pdfauthor={Maria Claudia Cabral da Silva},pdftitle={AV03: Construção do MVP de calibração}}
\setlength{\parindent}{0pt}\setlength{\parskip}{6pt}
\begin{document}
\begin{center}
{\sffamily\small UNISENAI \quad / \quad PROJETO APLICADO I \quad / \quad 2026.2}\\[1.2cm]
{\sffamily\bfseries\LARGE Calibração de dispositivos de controle}\\[.35cm]
{\large AV03: construção do MVP, parte 1}\\[.8cm]
Maria Claudia Cabral da Silva\\
Empresa parceira: Schulz S.A.\\
Docentes: Hugo Menezes Barra e Rodrigo Losina
\end{center}
\section*{Do requisito à consulta}
A AV1 identificou a necessidade de manter duas leituras por característica, preservar o sinal das tolerâncias e distinguir o resultado técnico da emissão documental. A AV2 definiu três tabelas SQLite. Nesta etapa, uma aplicação Python grava um evento sintético nesse mesmo esquema, consulta os cadastros relacionados e exibe o resultado por característica.

\begin{tabularx}{\linewidth}{@{}lX@{}}
\toprule\textbf{Objeto} & \textbf{Responsabilidade e vínculo com a AV2}\\\midrule
Tecnico & Representa o cadastro comum de técnico ou solicitante, com papel explícito e matrícula textual.\\
Equipamento & Identifica o dispositivo pela TAG, preservando código de origem e fabricante quando informados.\\
Medicao & Reúne as referências ao equipamento, técnico e solicitante; calcula média, erro e resultado de duas leituras.\\
Banco & Mapeia objetos às tabelas, executa CRUD e consulta as três entidades com JOIN.\\\bottomrule
\end{tabularx}

\section*{Uma medição, dois estados distintos}
Para cada característica, calcula-se a média das leituras e subtrai-se o nominal. O ensaio é aprovado quando o erro está entre os afastamentos inferior e superior, incluindo os extremos. Uma medição pode estar dentro da faixa e ainda aguardar revisão. O resultado do ensaio é derivado no domínio; o campo \texttt{status} da AV2 continua reservado ao fluxo documental.

Os dados de demonstração são fictícios. A classificação não substitui uma decisão metrológica homologada nem autoriza a emissão de um certificado.
\newpage
\section*{Execução e evidência}
Em uma pasta com os arquivos da entrega, execute \texttt{python 05\_main.py} com Python 3.10 ou superior. O programa cria uma base temporária, insere dois cadastros de papéis distintos, um equipamento e duas medições. A primeira medição é consultada, enviada para revisão e atualizada no banco. O relatório reúne TAG, técnico, solicitante, leituras e resultados por meio de uma consulta relacional.

\begin{center}
\small\begin{tabular}{@{}lrrrl@{}}
\toprule Parâmetro & Leituras (mm) & Média (mm) & Erro (mm) & Ensaio\\\midrule
Diâmetro A & 10,020 / 10,040 & 10,030 & +0,030 & APROVADO\\
Diâmetro B & 10,060 / 10,080 & 10,070 & +0,070 & REPROVADO\\\bottomrule
\end{tabular}
\end{center}
Nos dois casos, o nominal é 10,000 mm e os afastamentos são $\pm$0,050 mm. O diâmetro A permanece em revisão documental; o B permanece pendente. O terminal apresenta as leituras antes da média para permitir a conferência.

\section*{Integridade e testes}
As classes imutáveis validam entradas antes da gravação. O repositório usa parâmetros SQL, habilita as chaves estrangeiras a cada conexão e reverte a operação em caso de falha. O esquema original impede duplicar uma característica no evento e rejeita a troca dos papéis de técnico e solicitante na medição.

Execute \texttt{python -m unittest -v}. Os cinco testes passaram: foram conferidos os dois limites inclusivos, valores fora da tolerância, entradas inválidas, CRUD, JOIN, atualização dos valores derivados e rejeição de referências inválidas. A documentação Markdown contém o diagrama das classes e o detalhamento dos casos.

\section*{Arquivos e reprodução}
\texttt{03\_banco.sql} é o arquivo original da AV2; \texttt{04\_classes.py} implementa domínio e persistência; \texttt{05\_main.py} demonstra o fluxo; \texttt{06\_mvp\_parte1.md} documenta a solução. Para manter os dados após fechar o programa, execute \texttt{python 05\_main.py --banco demo.sqlite}.

{\small\textbf{Referências:} AV1, levantamento de requisitos de Maria Claudia Cabral da Silva; AV2, dicionário e esquema SQLite; roteiro AV03, itens 2 a 4.}
\end{document}
